% --- Archon physics formalization source begin ---
% archon:physics
% archon:covers PhyXMiniProblems/problem_phyx_mini_0360.lean
% archon:source-report reports/phyx_mini/problem_phyx_mini_0360.source.json

\chapter{Physics problem phyx\_mini\_0360}
\label{ch:PhyXMiniProblems_problem_phyx_mini_0360}

\paragraph{Problem source.}
\#\# Physical scenario

A large, 50.0-cm-diameter metal cylinder filled with air supports a 20.0 kg piston that can slide up and down without friction. The piston is 100.0 cm above the bottom when the temperature is 20\textasciicircum{}\textbackslash{}circ C. An 80.0 kg student then stands on the piston.

\#\# Question

After several minutes have elapsed, by how much has the piston been depressed?

\#\# Answer choices

- A: A: 3.8cm

- B: B: 3.0cm

- C: C: 5.0cm

- D: D: 4.0cm

\#\# Figure caption (auxiliary; use the image as primary evidence)

The image features two vertical cylinders filled with a fluid and covered with pistons. 

\#\#\# Left Cylinder:

- **Piston**: Labeled with mass \textbackslash{}( M = 20.0 \textbackslash{}, \textbackslash{}text\{kg\} \textbackslash{}).
- **Fluid below piston**: Associated with pressure \textbackslash{}( P\_1 \textbackslash{}).
- **Height of fluid**: \textbackslash{}( h\_1 \textbackslash{}).
- **Cross-sectional area**: Denoted as \textbackslash{}( \textbackslash{}text\{Area\} \textbackslash{}, A \textbackslash{}).

\#\#\# Right Cylinder:

- **Piston**: A stick figure is standing on it, labeled with mass \textbackslash{}( 80.0 \textbackslash{}, \textbackslash{}text\{kg\} \textbackslash{}).
- **Fluid below piston**: Associated with pressure \textbackslash{}( P\_2 \textbackslash{}).
- **Height of fluid**: \textbackslash{}( h\_2 \textbackslash{}).
- **Cross-sectional area**: Denoted as \textbackslash{}( \textbackslash{}text\{Area\} \textbackslash{}, A \textbackslash{}).

\#\#\# Additional Information:

- The atmospheric pressure is labeled as \textbackslash{}( P\_\{\textbackslash{}text\{atmos\}\} = 1 \textbackslash{}, \textbackslash{}text\{atm\} \textbackslash{}).

\#\# Dataset metadata

Ideal Gases and Kinetic Theory; Multi-Formula Reasoning, Physical Model Grounding Reasoning

\paragraph{Recorded answer/context.}
A: A: 3.8cm

\paragraph{Figure/image path.}
/root/proposal\_for\_physic/science-mango/phyx\_mini\_run/phyx\_data/test\_image/360.png

\paragraph{Formalization target.}
create a compiling Lean file with sorry bodies at `PhyXMiniProblems/problem\_phyx\_mini\_0360.lean`. The Lean declarations must preserve the physical quantities, dimensions or dimensional roles, figure labels, governing-law hypotheses, and final relation expressed by this problem.
Use Mathlib/Physlib names found through LeanExplore where available. If a physics API is missing, introduce faithful local abstractions rather than scalar placeholder aliases.

\begin{theorem}[Physics formalization target]
\label{thm:physics:phyx_mini_0360:target}
\lean{PhyXMiniProblems.ProblemPhyXMini0360.pistonDepression_rounds_to_three_point_eight_centimeters}
\uses{def:physics:phyx-mini-0360:phyxminiproblems-problemphyxmini0360-amountofsubstancescale, def:physics:phyx-mini-0360:phyxminiproblems-problemphyxmini0360-airpistonsetup, def:physics:phyx-mini-0360:phyxminiproblems-problemphyxmini0360-matchesproblemstatementandfigure, def:physics:phyx-mini-0360:phyxminiproblems-problemphyxmini0360-issealedandthermallyreequilibrated, def:physics:phyx-mini-0360:phyxminiproblems-problemphyxmini0360-usesstandardearthgravity, def:physics:phyx-mini-0360:phyxminiproblems-problemphyxmini0360-hasphysicalparameters, def:physics:phyx-mini-0360:phyxminiproblems-problemphyxmini0360-satisfiescylinderidealgasandequilibriumlaws, def:physics:phyx-mini-0360:phyxminiproblems-problemphyxmini0360-pistondepressionincentimeters, def:physics:phyx-mini-0360:phyxminiproblems-problemphyxmini0360-answerdepressionincentimeters}
The assigned autoformalize agent should translate this physics problem into Lean declarations in the covered file, with theorem and lemma proof bodies written as `by sorry`.
\end{theorem}
\begin{proof}
This is an autoformalization task, not a proof task. Produce statements that can later be proved by the physics prover without weakening the physical model.
\end{proof}

% --- Archon declaration topology begin ---
\section{Declaration topology}
\begin{definition}[Length Quantity]
  \label{def:physics:phyx-mini-0360:phyxminiproblems-problemphyxmini0360-lengthquantity}
  \lean{PhyXMiniProblems.ProblemPhyXMini0360.LengthQuantity}
  A physical length, independent of the chosen unit
\end{definition}
\begin{proof}
  This is the declared model definition of Length Quantity.
\end{proof}
\begin{definition}[Mass Quantity]
  \label{def:physics:phyx-mini-0360:phyxminiproblems-problemphyxmini0360-massquantity}
  \lean{PhyXMiniProblems.ProblemPhyXMini0360.MassQuantity}
  A physical mass, independent of the chosen unit
\end{definition}
\begin{proof}
  This is the declared model definition of Mass Quantity.
\end{proof}
\begin{definition}[Area Quantity]
  \label{def:physics:phyx-mini-0360:phyxminiproblems-problemphyxmini0360-areaquantity}
  \lean{PhyXMiniProblems.ProblemPhyXMini0360.AreaQuantity}
  A nonnegative physical cross-sectional area
\end{definition}
\begin{proof}
  This is the declared model definition of Area Quantity.
\end{proof}
\begin{definition}[Volume Quantity]
  \label{def:physics:phyx-mini-0360:phyxminiproblems-problemphyxmini0360-volumequantity}
  \lean{PhyXMiniProblems.ProblemPhyXMini0360.VolumeQuantity}
  A physical volume, with dimension length cubed
\end{definition}
\begin{proof}
  This is the declared model definition of Volume Quantity.
\end{proof}
\begin{definition}[Pressure Quantity]
  \label{def:physics:phyx-mini-0360:phyxminiproblems-problemphyxmini0360-pressurequantity}
  \lean{PhyXMiniProblems.ProblemPhyXMini0360.PressureQuantity}
  A physical pressure
\end{definition}
\begin{proof}
  This is the declared model definition of Pressure Quantity.
\end{proof}
\begin{definition}[Acceleration Quantity]
  \label{def:physics:phyx-mini-0360:phyxminiproblems-problemphyxmini0360-accelerationquantity}
  \lean{PhyXMiniProblems.ProblemPhyXMini0360.AccelerationQuantity}
  A physical acceleration, used for local gravitational acceleration
\end{definition}
\begin{proof}
  This is the declared model definition of Acceleration Quantity.
\end{proof}
\begin{definition}[Molar Gas Constant Quantity]
  \label{def:physics:phyx-mini-0360:phyxminiproblems-problemphyxmini0360-molargasconstantquantity}
  \lean{PhyXMiniProblems.ProblemPhyXMini0360.MolarGasConstantQuantity}
  The molar gas constant has energy-per-temperature dimension here. Its inverse amount-of-substance dimension is paired with the explicitly molar readout below because amount of substance is not a base dimension in Physlib
\end{definition}
\begin{proof}
  This is the declared model definition of Molar Gas Constant Quantity.
\end{proof}
\begin{definition}[Amount Of Substance Scale]
  \label{def:physics:phyx-mini-0360:phyxminiproblems-problemphyxmini0360-amountofsubstancescale}
  \lean{PhyXMiniProblems.ProblemPhyXMini0360.AmountOfSubstanceScale}
  An abstract amount-of-substance carrier together with its calibrated mole readout. This keeps the physical amount distinct from a bare real number
\end{definition}
\begin{proof}
  This is the declared model definition of Amount Of Substance Scale.
\end{proof}
\begin{definition}[length In Meters]
  \label{def:physics:phyx-mini-0360:phyxminiproblems-problemphyxmini0360-lengthinmeters}
  \lean{PhyXMiniProblems.ProblemPhyXMini0360.lengthInMeters}
  \uses{def:physics:phyx-mini-0360:phyxminiproblems-problemphyxmini0360-lengthquantity}
  Metre readout of a physical length
\end{definition}
\begin{proof}
  This is the declared model definition of length In Meters.
\end{proof}
\begin{definition}[mass In Kilograms]
  \label{def:physics:phyx-mini-0360:phyxminiproblems-problemphyxmini0360-massinkilograms}
  \lean{PhyXMiniProblems.ProblemPhyXMini0360.massInKilograms}
  \uses{def:physics:phyx-mini-0360:phyxminiproblems-problemphyxmini0360-massquantity}
  Kilogram readout of a physical mass
\end{definition}
\begin{proof}
  This is the declared model definition of mass In Kilograms.
\end{proof}
\begin{definition}[area In Square Meters]
  \label{def:physics:phyx-mini-0360:phyxminiproblems-problemphyxmini0360-areainsquaremeters}
  \lean{PhyXMiniProblems.ProblemPhyXMini0360.areaInSquareMeters}
  \uses{def:physics:phyx-mini-0360:phyxminiproblems-problemphyxmini0360-areaquantity}
  Square-metre readout of a physical area
\end{definition}
\begin{proof}
  This is the declared model definition of area In Square Meters.
\end{proof}
\begin{definition}[volume In Cubic Meters]
  \label{def:physics:phyx-mini-0360:phyxminiproblems-problemphyxmini0360-volumeincubicmeters}
  \lean{PhyXMiniProblems.ProblemPhyXMini0360.volumeInCubicMeters}
  \uses{def:physics:phyx-mini-0360:phyxminiproblems-problemphyxmini0360-volumequantity}
  Cubic-metre readout of a physical volume
\end{definition}
\begin{proof}
  This is the declared model definition of volume In Cubic Meters.
\end{proof}
\begin{definition}[pressure In Pascals]
  \label{def:physics:phyx-mini-0360:phyxminiproblems-problemphyxmini0360-pressureinpascals}
  \lean{PhyXMiniProblems.ProblemPhyXMini0360.pressureInPascals}
  \uses{def:physics:phyx-mini-0360:phyxminiproblems-problemphyxmini0360-pressurequantity}
  Pascal readout of a physical pressure
\end{definition}
\begin{proof}
  This is the declared model definition of pressure In Pascals.
\end{proof}
\begin{definition}[acceleration In Meters Per Second Squared]
  \label{def:physics:phyx-mini-0360:phyxminiproblems-problemphyxmini0360-accelerationinmeterspersecondsquared}
  \lean{PhyXMiniProblems.ProblemPhyXMini0360.accelerationInMetersPerSecondSquared}
  \uses{def:physics:phyx-mini-0360:phyxminiproblems-problemphyxmini0360-accelerationquantity}
  Metres-per-second-squared readout of a physical acceleration
\end{definition}
\begin{proof}
  This is the declared model definition of acceleration In Meters Per Second Squared.
\end{proof}
\begin{definition}[temperature In Kelvin]
  \label{def:physics:phyx-mini-0360:phyxminiproblems-problemphyxmini0360-temperatureinkelvin}
  \lean{PhyXMiniProblems.ProblemPhyXMini0360.temperatureInKelvin}
  The scenario fixes Physlib's arbitrary absolute-temperature scale to kelvin, so Temperature.toReal is the kelvin readout in this setup
\end{definition}
\begin{proof}
  This is the declared model definition of temperature In Kelvin.
\end{proof}
\begin{definition}[temperature In Celsius]
  \label{def:physics:phyx-mini-0360:phyxminiproblems-problemphyxmini0360-temperatureincelsius}
  \lean{PhyXMiniProblems.ProblemPhyXMini0360.temperatureInCelsius}
  \uses{def:physics:phyx-mini-0360:phyxminiproblems-problemphyxmini0360-temperatureinkelvin}
  Celsius readout derived from an absolute kelvin readout
\end{definition}
\begin{proof}
  This is the declared model definition of temperature In Celsius.
\end{proof}
\begin{definition}[molar Gas Constant In SI]
  \label{def:physics:phyx-mini-0360:phyxminiproblems-problemphyxmini0360-molargasconstantinsi}
  \lean{PhyXMiniProblems.ProblemPhyXMini0360.molarGasConstantInSI}
  \uses{def:physics:phyx-mini-0360:phyxminiproblems-problemphyxmini0360-molargasconstantquantity}
  SI readout of the molar gas constant
\end{definition}
\begin{proof}
  This is the declared model definition of molar Gas Constant In SI.
\end{proof}
\begin{definition}[Equilibrium State]
  \label{def:physics:phyx-mini-0360:phyxminiproblems-problemphyxmini0360-equilibriumstate}
  \lean{PhyXMiniProblems.ProblemPhyXMini0360.EquilibriumState}
  The two mechanical-equilibrium states shown side by side in the figure
\end{definition}
\begin{proof}
  This is the declared model definition of Equilibrium State.
\end{proof}
\begin{definition}[Figure Pressure Label]
  \label{def:physics:phyx-mini-0360:phyxminiproblems-problemphyxmini0360-figurepressurelabel}
  \lean{PhyXMiniProblems.ProblemPhyXMini0360.FigurePressureLabel}
  Pressure symbols printed in the primary figure
\end{definition}
\begin{proof}
  This is the declared model definition of Figure Pressure Label.
\end{proof}
\begin{definition}[Figure Height Label]
  \label{def:physics:phyx-mini-0360:phyxminiproblems-problemphyxmini0360-figureheightlabel}
  \lean{PhyXMiniProblems.ProblemPhyXMini0360.FigureHeightLabel}
  Height symbols printed in the primary figure
\end{definition}
\begin{proof}
  This is the declared model definition of Figure Height Label.
\end{proof}
\begin{definition}[Figure Area Label]
  \label{def:physics:phyx-mini-0360:phyxminiproblems-problemphyxmini0360-figurearealabel}
  \lean{PhyXMiniProblems.ProblemPhyXMini0360.FigureAreaLabel}
  Cross-sectional-area symbol printed in both figure panels
\end{definition}
\begin{proof}
  This is the declared model definition of Figure Area Label.
\end{proof}
\begin{definition}[Cylinder Shape]
  \label{def:physics:phyx-mini-0360:phyxminiproblems-problemphyxmini0360-cylindershape}
  \lean{PhyXMiniProblems.ProblemPhyXMini0360.CylinderShape}
  Shape and orientation of the container
\end{definition}
\begin{proof}
  This is the declared model definition of Cylinder Shape.
\end{proof}
\begin{definition}[Cylinder Material]
  \label{def:physics:phyx-mini-0360:phyxminiproblems-problemphyxmini0360-cylindermaterial}
  \lean{PhyXMiniProblems.ProblemPhyXMini0360.CylinderMaterial}
  Material named for the cylinder
\end{definition}
\begin{proof}
  This is the declared model definition of Cylinder Material.
\end{proof}
\begin{definition}[Gas Species]
  \label{def:physics:phyx-mini-0360:phyxminiproblems-problemphyxmini0360-gasspecies}
  \lean{PhyXMiniProblems.ProblemPhyXMini0360.GasSpecies}
  Gas species named in the problem
\end{definition}
\begin{proof}
  This is the declared model definition of Gas Species.
\end{proof}
\begin{definition}[Gas Model]
  \label{def:physics:phyx-mini-0360:phyxminiproblems-problemphyxmini0360-gasmodel}
  \lean{PhyXMiniProblems.ProblemPhyXMini0360.GasModel}
  Equation-of-state model used for the trapped air
\end{definition}
\begin{proof}
  This is the declared model definition of Gas Model.
\end{proof}
\begin{definition}[Piston Motion]
  \label{def:physics:phyx-mini-0360:phyxminiproblems-problemphyxmini0360-pistonmotion}
  \lean{PhyXMiniProblems.ProblemPhyXMini0360.PistonMotion}
  Constraint on the piston's allowed motion
\end{definition}
\begin{proof}
  This is the declared model definition of Piston Motion.
\end{proof}
\begin{definition}[Above Piston Medium]
  \label{def:physics:phyx-mini-0360:phyxminiproblems-problemphyxmini0360-abovepistonmedium}
  \lean{PhyXMiniProblems.ProblemPhyXMini0360.AbovePistonMedium}
  Medium exerting pressure on the top face of the piston
\end{definition}
\begin{proof}
  This is the declared model definition of Above Piston Medium.
\end{proof}
\begin{definition}[Additional Load]
  \label{def:physics:phyx-mini-0360:phyxminiproblems-problemphyxmini0360-additionalload}
  \lean{PhyXMiniProblems.ProblemPhyXMini0360.AdditionalLoad}
  How the second load is applied
\end{definition}
\begin{proof}
  This is the declared model definition of Additional Load.
\end{proof}
\begin{definition}[Waiting Protocol]
  \label{def:physics:phyx-mini-0360:phyxminiproblems-problemphyxmini0360-waitingprotocol}
  \lean{PhyXMiniProblems.ProblemPhyXMini0360.WaitingProtocol}
  Waiting condition described before the requested final readout
\end{definition}
\begin{proof}
  This is the declared model definition of Waiting Protocol.
\end{proof}
\begin{definition}[Air Piston Setup]
  \label{def:physics:phyx-mini-0360:phyxminiproblems-problemphyxmini0360-airpistonsetup}
  \lean{PhyXMiniProblems.ProblemPhyXMini0360.AirPistonSetup}
  \uses{def:physics:phyx-mini-0360:phyxminiproblems-problemphyxmini0360-lengthquantity, def:physics:phyx-mini-0360:phyxminiproblems-problemphyxmini0360-massquantity, def:physics:phyx-mini-0360:phyxminiproblems-problemphyxmini0360-areaquantity, def:physics:phyx-mini-0360:phyxminiproblems-problemphyxmini0360-volumequantity, def:physics:phyx-mini-0360:phyxminiproblems-problemphyxmini0360-pressurequantity, def:physics:phyx-mini-0360:phyxminiproblems-problemphyxmini0360-accelerationquantity, def:physics:phyx-mini-0360:phyxminiproblems-problemphyxmini0360-molargasconstantquantity, def:physics:phyx-mini-0360:phyxminiproblems-problemphyxmini0360-amountofsubstancescale, def:physics:phyx-mini-0360:phyxminiproblems-problemphyxmini0360-equilibriumstate, def:physics:phyx-mini-0360:phyxminiproblems-problemphyxmini0360-figurepressurelabel, def:physics:phyx-mini-0360:phyxminiproblems-problemphyxmini0360-figureheightlabel, def:physics:phyx-mini-0360:phyxminiproblems-problemphyxmini0360-figurearealabel, def:physics:phyx-mini-0360:phyxminiproblems-problemphyxmini0360-cylindershape, def:physics:phyx-mini-0360:phyxminiproblems-problemphyxmini0360-cylindermaterial, def:physics:phyx-mini-0360:phyxminiproblems-problemphyxmini0360-gasspecies, def:physics:phyx-mini-0360:phyxminiproblems-problemphyxmini0360-gasmodel, def:physics:phyx-mini-0360:phyxminiproblems-problemphyxmini0360-pistonmotion, def:physics:phyx-mini-0360:phyxminiproblems-problemphyxmini0360-abovepistonmedium, def:physics:phyx-mini-0360:phyxminiproblems-problemphyxmini0360-additionalload, def:physics:phyx-mini-0360:phyxminiproblems-problemphyxmini0360-waitingprotocol}
  The physical apparatus and its two independent equilibrium states. The pressure and height functions are the figure's P₁, P₂, h₁, and h₂. In particular, the final height is an unconstrained field and is not defined from the displayed 3.8 cm answer
\end{definition}
\begin{proof}
  This is the declared model definition of Air Piston Setup.
\end{proof}
\begin{definition}[supported Mass In Kilograms]
  \label{def:physics:phyx-mini-0360:phyxminiproblems-problemphyxmini0360-supportedmassinkilograms}
  \lean{PhyXMiniProblems.ProblemPhyXMini0360.supportedMassInKilograms}
  \uses{def:physics:phyx-mini-0360:phyxminiproblems-problemphyxmini0360-amountofsubstancescale, def:physics:phyx-mini-0360:phyxminiproblems-problemphyxmini0360-massinkilograms, def:physics:phyx-mini-0360:phyxminiproblems-problemphyxmini0360-equilibriumstate, def:physics:phyx-mini-0360:phyxminiproblems-problemphyxmini0360-airpistonsetup}
  Total load mass supported by the pressure difference in each state. This only records which bodies stand on the piston; it does not constrain a height
\end{definition}
\begin{proof}
  This is the declared model definition of supported Mass In Kilograms.
\end{proof}
\begin{definition}[Matches Problem Statement And Figure]
  \label{def:physics:phyx-mini-0360:phyxminiproblems-problemphyxmini0360-matchesproblemstatementandfigure}
  \lean{PhyXMiniProblems.ProblemPhyXMini0360.MatchesProblemStatementAndFigure}
  \uses{def:physics:phyx-mini-0360:phyxminiproblems-problemphyxmini0360-amountofsubstancescale, def:physics:phyx-mini-0360:phyxminiproblems-problemphyxmini0360-lengthinmeters, def:physics:phyx-mini-0360:phyxminiproblems-problemphyxmini0360-massinkilograms, def:physics:phyx-mini-0360:phyxminiproblems-problemphyxmini0360-temperatureincelsius, def:physics:phyx-mini-0360:phyxminiproblems-problemphyxmini0360-airpistonsetup}
  Qualitative setup, primary-figure label assignments, and numerical readouts given in the source. No final height or depression value occurs here
\end{definition}
\begin{proof}
  This is the declared model definition of Matches Problem Statement And Figure.
\end{proof}
\begin{definition}[Is Sealed And Thermally Reequilibrated]
  \label{def:physics:phyx-mini-0360:phyxminiproblems-problemphyxmini0360-issealedandthermallyreequilibrated}
  \lean{PhyXMiniProblems.ProblemPhyXMini0360.IsSealedAndThermallyReequilibrated}
  \uses{def:physics:phyx-mini-0360:phyxminiproblems-problemphyxmini0360-amountofsubstancescale, def:physics:phyx-mini-0360:phyxminiproblems-problemphyxmini0360-airpistonsetup}
  The piston is observed after the sealed gas has thermally re-equilibrated with the unchanged ambient. These are operating conditions, not a final-height readout
\end{definition}
\begin{proof}
  This is the declared model definition of Is Sealed And Thermally Reequilibrated.
\end{proof}
\begin{definition}[Uses Standard Earth Gravity]
  \label{def:physics:phyx-mini-0360:phyxminiproblems-problemphyxmini0360-usesstandardearthgravity}
  \lean{PhyXMiniProblems.ProblemPhyXMini0360.UsesStandardEarthGravity}
  \uses{def:physics:phyx-mini-0360:phyxminiproblems-problemphyxmini0360-amountofsubstancescale, def:physics:phyx-mini-0360:phyxminiproblems-problemphyxmini0360-accelerationinmeterspersecondsquared, def:physics:phyx-mini-0360:phyxminiproblems-problemphyxmini0360-airpistonsetup}
  The standard near-surface Earth acceleration used for the numerical answer. This environmental calibration is independent of the unknown displacement
\end{definition}
\begin{proof}
  This is the declared model definition of Uses Standard Earth Gravity.
\end{proof}
\begin{definition}[Has Physical Parameters]
  \label{def:physics:phyx-mini-0360:phyxminiproblems-problemphyxmini0360-hasphysicalparameters}
  \lean{PhyXMiniProblems.ProblemPhyXMini0360.HasPhysicalParameters}
  \uses{def:physics:phyx-mini-0360:phyxminiproblems-problemphyxmini0360-amountofsubstancescale, def:physics:phyx-mini-0360:phyxminiproblems-problemphyxmini0360-lengthinmeters, def:physics:phyx-mini-0360:phyxminiproblems-problemphyxmini0360-massinkilograms, def:physics:phyx-mini-0360:phyxminiproblems-problemphyxmini0360-areainsquaremeters, def:physics:phyx-mini-0360:phyxminiproblems-problemphyxmini0360-volumeincubicmeters, def:physics:phyx-mini-0360:phyxminiproblems-problemphyxmini0360-pressureinpascals, def:physics:phyx-mini-0360:phyxminiproblems-problemphyxmini0360-accelerationinmeterspersecondsquared, def:physics:phyx-mini-0360:phyxminiproblems-problemphyxmini0360-temperatureinkelvin, def:physics:phyx-mini-0360:phyxminiproblems-problemphyxmini0360-molargasconstantinsi, def:physics:phyx-mini-0360:phyxminiproblems-problemphyxmini0360-airpistonsetup}
  Positivity and nondegeneracy conditions selecting the physical branch
\end{definition}
\begin{proof}
  This is the declared model definition of Has Physical Parameters.
\end{proof}
\begin{definition}[Satisfies Cylinder Ideal Gas And Equilibrium Laws]
  \label{def:physics:phyx-mini-0360:phyxminiproblems-problemphyxmini0360-satisfiescylinderidealgasandequilibriumlaws}
  \lean{PhyXMiniProblems.ProblemPhyXMini0360.SatisfiesCylinderIdealGasAndEquilibriumLaws}
  \uses{def:physics:phyx-mini-0360:phyxminiproblems-problemphyxmini0360-amountofsubstancescale, def:physics:phyx-mini-0360:phyxminiproblems-problemphyxmini0360-lengthinmeters, def:physics:phyx-mini-0360:phyxminiproblems-problemphyxmini0360-areainsquaremeters, def:physics:phyx-mini-0360:phyxminiproblems-problemphyxmini0360-volumeincubicmeters, def:physics:phyx-mini-0360:phyxminiproblems-problemphyxmini0360-pressureinpascals, def:physics:phyx-mini-0360:phyxminiproblems-problemphyxmini0360-accelerationinmeterspersecondsquared, def:physics:phyx-mini-0360:phyxminiproblems-problemphyxmini0360-temperatureinkelvin, def:physics:phyx-mini-0360:phyxminiproblems-problemphyxmini0360-molargasconstantinsi, def:physics:phyx-mini-0360:phyxminiproblems-problemphyxmini0360-airpistonsetup, def:physics:phyx-mini-0360:phyxminiproblems-problemphyxmini0360-supportedmassinkilograms}
  The governing circular-cylinder geometry, ideal-gas equation, and static vertical force balance. The force law is (P\_gas - P\_atmos) A = m\_supported g. These relations apply uniformly to the two independent states and contain no numerical value for h₂ or for the requested depression
\end{definition}
\begin{proof}
  This is the declared model definition of Satisfies Cylinder Ideal Gas And Equilibrium Laws.
\end{proof}
\begin{definition}[piston Depression In Centimeters]
  \label{def:physics:phyx-mini-0360:phyxminiproblems-problemphyxmini0360-pistondepressionincentimeters}
  \lean{PhyXMiniProblems.ProblemPhyXMini0360.pistonDepressionInCentimeters}
  \uses{def:physics:phyx-mini-0360:phyxminiproblems-problemphyxmini0360-amountofsubstancescale, def:physics:phyx-mini-0360:phyxminiproblems-problemphyxmini0360-lengthinmeters, def:physics:phyx-mini-0360:phyxminiproblems-problemphyxmini0360-airpistonsetup}
  The signed downward displacement in centimetres. This definition is only the geometric difference 100 (h₁ - h₂) and contains no answer-choice value
\end{definition}
\begin{proof}
  This is the declared model definition of piston Depression In Centimeters.
\end{proof}
\begin{definition}[Answer Choice]
  \label{def:physics:phyx-mini-0360:phyxminiproblems-problemphyxmini0360-answerchoice}
  \lean{PhyXMiniProblems.ProblemPhyXMini0360.AnswerChoice}
  Labels of the four answers printed in the source
\end{definition}
\begin{proof}
  This is the declared model definition of Answer Choice.
\end{proof}
\begin{definition}[answer Depression In Centimeters]
  \label{def:physics:phyx-mini-0360:phyxminiproblems-problemphyxmini0360-answerdepressionincentimeters}
  \lean{PhyXMiniProblems.ProblemPhyXMini0360.answerDepressionInCentimeters}
  \uses{def:physics:phyx-mini-0360:phyxminiproblems-problemphyxmini0360-answerchoice}
  Centimetre value printed beside each answer label
\end{definition}
\begin{proof}
  This is the declared model definition of answer Depression In Centimeters.
\end{proof}
\begin{definition}[recorded Answer Choice]
  \label{def:physics:phyx-mini-0360:phyxminiproblems-problemphyxmini0360-recordedanswerchoice}
  \lean{PhyXMiniProblems.ProblemPhyXMini0360.recordedAnswerChoice}
  \uses{def:physics:phyx-mini-0360:phyxminiproblems-problemphyxmini0360-answerchoice}
  The answer label recorded in the supplied dataset
\end{definition}
\begin{proof}
  This is the declared model definition of recorded Answer Choice.
\end{proof}
% --- Archon declaration topology end ---
% --- Archon physics formalization source end ---
